This module provides the finite-dimensional analytic input used to verify:
\begin{enumerate}[label=(G\arabic*)]
\item the \emph{uniform convexity} hypothesis in the Wilson-class setting for the given compact gauge group $G$ (on a small-curvature chart domain), with explicit constants in principle; and
\item a canonical extracted-irrelevant scaling gap $\Delta_{\min}=2$ once all operators of engineering dimension $\le 4$ are extracted.
\end{enumerate}
With these finite-dimensional checks in hand (uniform convexity and an extracted-irrelevant scaling gap), the remaining components are:
(a) implement the good-chart event and (b) fix the extraction bookkeeping.

\paragraph{Scope note.}
The convexity statement proved here is a \emph{local chart-domain} theorem for one coarse block in the chart-truncated Lane~C map.
It is logically separate from the repaired weak-window package in Appendix~\ref{app:route1}; Theorem~\ref{thm:weak_window_sufficiency}
there certifies the downstream step-scaling / Route~1 intake branch without any additional convexity input beyond that appendix-level package.

We fix one dyadic RG step $j\to j+1$ and work on a single unit coarse block $\mathcal Q$ (a $2\times 2\times 2\times 2$ hypercube in $j$-block units).
Let $E(\mathcal Q)$ be its internal oriented edges and $P(\mathcal Q)$ its oriented plaquettes.
Let $d_G:=\dim G$.

For $A=(A_e)_{e\in E(\mathcal Q)}$ with $A_e\in\mathfrak g\cong\mathbb R^{d_G}$, define the \emph{abelianized plaquette flux} (discrete curl)
\[
(D A)_p\ :=\ \sum_{e\in\partial p} \sigma_{p,e}\,A_e,
\qquad \sigma_{p,e}\in\{-1,0,1\}.
\]
After imposing tree gauge (fixing a maximal tree of edges to $A_e=0$), $D$ becomes a fixed linear map
\[
D_{\mathrm{red}}:\mathbb R^{n_{\mathrm{loc}}}\to \mathbb R^{n_{\mathrm{plaq}}},
\]
where $n_{\mathrm{loc}}$ and $n_{\mathrm{plaq}}$ are fixed integers \emph{per Lie component}.
For the canonical $2^4$ block and canonical tree gauge, one has $n_{\mathrm{loc}}^{(1)}=17$ and $n_{\mathrm{plaq}}^{(1)}=24$ for a single Lie component,
hence $n_{\mathrm{loc}}=d_G\,n_{\mathrm{loc}}^{(1)}$ and $n_{\mathrm{plaq}}=d_G\,n_{\mathrm{plaq}}^{(1)}$ in total.

\begin{definition}[Curl coercivity constant]\label{def:lambda_star}
Define
\[
\lambda_\ast\ :=\ \lambda_{\min}\big(D_{\mathrm{red}}^\ast D_{\mathrm{red}}\big)\ >\ 0.
\]
Then for all $v\in\mathbb R^{n_{\mathrm{loc}}}$,
\begin{equation}\label{eq:curl_coercive}
\|D_{\mathrm{red}}v\|^2\ \ge\ \lambda_\ast\,\|v\|^2.
\end{equation}
\end{definition}

\begin{lemma}[Explicit lower bound $\lambda_\ast>1/4$ for the $2^4$ block]\label{lem:lambda_lower}
For the $2\times2\times2\times2$ coarse block with the canonical tree gauge, the reduced integer matrix
$Q:=D_{\mathrm{red}}^\ast D_{\mathrm{red}}$ (for \emph{one} Lie component) has smallest eigenvalue $\lambda_\ast$ satisfying
\[
\lambda_\ast\ >\ \frac14.
\]
Consequently the same lower bound holds for the full $G$ operator $D_{\mathrm{red}}^\ast D_{\mathrm{red}}$ since it is block-diagonal over Lie components.
\end{lemma}

\begin{proof}
$Q$ is a fixed $17\times 17$ symmetric integer matrix (for one Lie component).
Appendix~\ref{app:lambda_star_Q_to_P} records the explicit matrix $Q$ (in the canonical tree gauge) and provides an explicit algebraic verification of the exact factorization
of its characteristic polynomial $\chi_Q(x)=\det(xI-Q)$:
\[
\chi_Q(x)=(x-6)(x-1)^3(x^3-8x^2+16x-6)\,P(x),
\]
where $P$ is the degree-$10$ polynomial displayed in \eqref{eq:P_lambda_star}.
The remaining factor $(x-6)(x-1)^3(x^3-8x^2+16x-6)$ has no real roots in $[0,\tfrac14]$ (Appendix~\ref{app:lambda_star_Q_to_P}).
Therefore $\lambda_\ast=\lambda_{\min}(Q)>\tfrac14$ follows once $P$ has no real roots in $[0,\tfrac14]$.
A complete Sturm--Habicht (subresultant) root-isolation argument for this claim is included in Appendix~\ref{app:lambda_star_certificate}.
\end{proof}

\paragraph{Remark.}
For the convexity argument, only $\lambda_\ast>0$ is needed. Lemma~\ref{lem:lambda_lower} gives a concrete quantitative margin and,
crucially, $\lambda_\ast$ is \emph{independent of $G$} (and hence of the rank parameter $n$).

\subsection{Chart coordinates and BCH control}

Fix a chart radius $\delta\in(0,\pi/4]$. Parameterize each unfixed link by $U_e=\exp(A_e)$ with $\|A_e\|\le\delta$.
For each plaquette $p$, define the plaquette holonomy
\[
U_p(A)\ :=\ \prod_{e\in\partial p} U_e^{\sigma_{p,e}}\ \in G,
\qquad
F_p(A)\ :=\ \log\big(U_p(A)\big)\in\mathfrak g,
\]
using the principal logarithm on the ball $\|F\|\le \pi$.

\begin{lemma}[BCH decomposition and derivative bounds]\label{lem:BCH_bounds}
There exist constants $C_2(\delta;G),C_1(\delta;G),C_0(\delta;G)$ such that for all $A$ with $\|A_e\|\le \delta$:
\begin{enumerate}[label=(B\arabic*)]
\item (\textbf{Quadratic nonlinearity}) $F(A)$ decomposes as
\begin{equation}\label{eq:F_decomp}
F(A)\ =\ D_{\mathrm{red}}A\ +\ N(A),
\end{equation}
where $N(A)$ satisfies $\|N(A)\|\le C_2(\delta;G)\,\|A\|^2$.
\item (\textbf{Jacobian perturbation}) The Jacobian satisfies $\|\nabla N(A)\|\le C_1(\delta;G)\,\|A\|$.
\item (\textbf{Second derivative}) The second derivative satisfies $\|\nabla^2 N(A)\|\le C_0(\delta;G)$.
\end{enumerate}
All norms are Euclidean/operator norms on the finite-dimensional local coordinate space.
\end{lemma}

\paragraph{Comment.}
These are standard BCH remainder bounds on a compact Lie algebra.
For fixed $G$ one may take $C_i(\delta;G)$ polynomial in $\delta$ for $\delta\le 1$.
For the purposes of the multiscale construction, only finiteness and scale-independence (uniform in the dyadic scale $j$) are needed.

\subsection{Wilson potential and Haar Jacobian}

For the (possibly anisotropic) Wilson-class plaquette weight we use the local potential
\[
V_{\beta,\alpha}(X)\ :=\ \alpha\beta\Big(1-\frac{1}{\dim\rho}\mathrm{Re\,Tr}_{\rho}(\exp X)\Big),
\qquad X\in\mathfrak g,\qquad \alpha\in[\alpha_{\min},1].
\]
It is analytic and Ad-invariant. Its Hessian at $X=0$ is positive definite.

\begin{lemma}[Wilson potential Hessian lower bound]\label{lem:radial_hess}
Fix $r_\ast\in(0,1]$. There exist constants $m(r_\ast;G)>0$ and $M(r_\ast;G)<\infty$ such that for all $X\in\mathfrak g$ with $\|X\|\le r_\ast$
and all $\alpha\in[\alpha_{\min},1]$,
\begin{equation}\label{eq:V_hess_lb}
\nabla_X^2 V_{\beta,\alpha}(X)\ \succeq\ \alpha\beta\,m(r_\ast;G)\,I,
\end{equation}
and
\begin{equation}\label{eq:V_grad_bd}
\|\nabla_X V_{\beta,\alpha}(X)\|\ \le\ \alpha\beta\,M(r_\ast;G)\,\|X\|.
\end{equation}
\end{lemma}

\begin{proof}[Proof (explicit constants in principle)]
Let $\rho$ denote the fixed finite-dimensional unitary Wilson representation from \eqref{eq:families}.
Set
\[
W(X):=1-\frac{1}{\dim\rho}\Re\Tr_{\rho}(e^{X}),\qquad V_{\beta,\alpha}(X)=\alpha\beta\,W(X).
\]
The map $W:\mathfrak g\to\mathbb R$ is real-analytic and Ad-invariant, with $W(0)=0$ and $\nabla W(0)=0$.
Define the (positive) Dynkin index $I_{\rho}>0$ relative to the fixed inner product by
\[
-\Tr_{\rho}(\rho_{\ast}(X)\,\rho_{\ast}(Y))\ =\ I_{\rho}\,\langle X,Y\rangle\qquad (X,Y\in\mathfrak g),
\]
where $\rho_{\ast}$ is the differential of $\rho$.
A Taylor expansion gives
$\Re\Tr_{\rho}(e^{X})=\dim\rho+\tfrac12\Re\Tr_{\rho}(X^{2})+O(\|X\|^{3})$ and $\Re\Tr_{\rho}(X)=0$ since $X$ is skew-adjoint.
Thus
\[
W(X)=\frac{1}{2\dim\rho}\bigl(-\Re\Tr_{\rho}(X^{2})\bigr)+O(\|X\|^{3})
=\frac{I_{\rho}}{2\dim\rho}\,\|X\|^{2}+O(\|X\|^{3}),
\]
so $\nabla^{2}W(0)=\frac{I_{\rho}}{\dim\rho}\,I$.
By compactness of $\{\|X\|\le r_\ast\}$ and continuity of $\nabla^{2}W$, the quantity
\[ m(r_{\ast};G):=\inf_{\|X\|\le r_\ast}\lambda_{\min}(\nabla^{2}W(X)) \]
exists and is strictly positive.
Multiplying by $\alpha\beta$ yields \eqref{eq:V_hess_lb}.
Similarly, since $\nabla W(0)=0$ and $\nabla W$ is Lipschitz on $\{\|X\|\le r_\ast\}$, there exists finite $M(r_\ast;G)$ such that
$\|\nabla W(X)\|\le M(r_\ast;G)\,\|X\|$, and multiplying by $\alpha\beta$ yields \eqref{eq:V_grad_bd}.
\end{proof}

On a unit coarse block $\mathcal Q$, define the local negative log-density
\begin{equation}\label{eq:H_def}
H_{j+1}(A)\ :=\ \sum_{p\in P(\mathcal Q)} V_{\beta,\alpha_p}\big(F_p(A)\big)\ -\ \sum_{e\in E_{\mathrm{red}}(\mathcal Q)}\log J_{\mathrm{Haar}}(A_e),
\end{equation}
where $\alpha_p\in[\alpha_{\min},1]$ are the local exponents (anisotropy/thickening parameters),
and $J_{\mathrm{Haar}}(A)$ is the Haar Jacobian of $G$ in exponential coordinates.
A convenient closed form is
\[
J_{\mathrm{Haar}}(A)\ =\ \det\!\Big(\frac{1-e^{-\mathrm{ad}_A}}{\mathrm{ad}_A}\Big),
\]
and for $A$ conjugate into a Cartan subalgebra this equals the familiar Weyl product
$J_{\mathrm{Haar}}(A)=\prod_{\alpha>0}\big(\frac{\sin(\alpha(A)/2)}{\alpha(A)/2}\big)^2$.

\begin{lemma}[Haar Jacobian convexity]\label{lem:haar_convex}
Let $\mathfrak t\subset\mathfrak g$ be a fixed Cartan subalgebra and $\Phi\subset\mathfrak t^{\ast}$ the corresponding root system.
Set
\[
L_G:=\max_{\alpha\in\Phi}\|\alpha\|,
\qquad
\delta_{\mathrm{Haar}}(G):=\frac{\pi}{L_G},
\]
where $\|\alpha\|$ is the operator norm of the real linear functional $\alpha:\mathfrak t\to\mathbb R$ induced by the fixed inner product on $\mathfrak g$.
Then for all $A\in\mathfrak g$ with $\|A\|\le \delta_{\mathrm{Haar}}(G)$ one has $J_{\mathrm{Haar}}(A)>0$ and the function $-\log J_{\mathrm{Haar}}(A)$ is convex:
\[
\nabla^2\big(-\log J_{\mathrm{Haar}}(A)\big)\ \succeq\ 0.
\]
\end{lemma}

\begin{proof}
We first work on $\mathfrak t$. For $H\in\mathfrak t$ with $\|H\|\le \delta_{\mathrm{Haar}}(G)$ we have $|\alpha(H)|\le \|\alpha\|\,\|H\|\le \pi$ for all $\alpha\in\Phi$.
Using the Weyl product formula for the Jacobian of the exponential map,
\[
J_{\mathrm{Haar}}(H)=\prod_{\alpha>0}\Big(\frac{\sin(\alpha(H)/2)}{\alpha(H)/2}\Big)^2,
\]
and since $t/\sin t>0$ for $t\in(-\pi/2,\pi/2)$ (hence for $t=\alpha(H)/2$ here), we may write
\[
-\log J_{\mathrm{Haar}}(H)=\sum_{\alpha>0}\phi\!\Big(\frac{\alpha(H)}{2}\Big),
\qquad
\phi(t):=2\log\Big(\frac{t}{\sin t}\Big).
\]
A direct computation gives, for $t\neq 0$,
\[
\phi''(t)=2\Big(\csc^2 t-\frac{1}{t^2}\Big)=2\Big(\frac{1}{\sin^2 t}-\frac{1}{t^2}\Big)\ \ge\ 0,
\]
because $|\sin t|\le |t|$ for all real $t$.
Defining $\phi''(0):=\lim_{t\to 0}\phi''(t)=\tfrac{2}{3}$, we have $\phi''(t)\ge 0$ on $(-\pi/2,\pi/2)$, hence $\phi$ is convex there.
Since each map $H\mapsto \alpha(H)/2$ is linear on $\mathfrak t$, each summand $H\mapsto \phi(\alpha(H)/2)$ is convex, and so is their finite sum.
Thus $-\log J_{\mathrm{Haar}}$ is convex on $\{H\in\mathfrak t:\|H\|\le \delta_{\mathrm{Haar}}(G)\}$.

To extend to $\mathfrak g$, note that $J_{\mathrm{Haar}}$ is $\Ad$-invariant and real-analytic, hence so is $f(A):=-\log J_{\mathrm{Haar}}(A)$.
For any $g\in G$, the map $A\mapsto \Ad_g A$ is an isometry for the fixed bi-invariant inner product, and by $\Ad$-invariance of $f$,
$f(\Ad_g A)=f(A)$.
Differentiating twice yields
\[
\nabla^2 f(\Ad_g A)=\Ad_g\circ (\nabla^2 f(A))\circ \Ad_{g^{-1}},
\]
so positive semidefiniteness of the Hessian is preserved along $\Ad$-orbits.
Every $A\in\mathfrak g$ is conjugate to some $H\in\mathfrak t$ with $\|H\|=\|A\|$.
Hence if $\|A\|\le \delta_{\mathrm{Haar}}(G)$ we may choose such an $H$ in the above Cartan neighborhood and conclude $\nabla^2 f(A)\succeq 0$.
\end{proof}


\subsection{Uniform convexity on the chart domain}

\begin{theorem}[Uniform convexity on the chart domain]\label{thm:uniform_convexity}
Fix $\delta\le \min\{\pi/4,\delta_{\mathrm{Haar}}(G)\}$ and define
\[
r_\ast:=\sup_{\|A\|_{\infty}\le \delta}\ \max_{p\in P(\mathcal Q)} \|F_p(A)\|.
\]
Assume the BCH derivative bounds of Lemma~\ref{lem:BCH_bounds} and define
\[
\varepsilon_1(\delta):=\sup_{\|A\|_\infty\le \delta}\|\nabla N(A)\|,
\qquad
\varepsilon_2(\delta):=\sup_{\|A\|_\infty\le \delta}\|\nabla^2 N(A)\|.
\]
Assume the smallness conditions
\begin{align}
\varepsilon_1(\delta)&\le \tfrac12\sqrt{\lambda_\ast},\label{eq:small1}\\
M(r_\ast;G)\,\varepsilon_2(\delta)\,r_\ast&\le \tfrac{\alpha_{\min}}{8}\,m(r_\ast;G)\,\lambda_\ast.\label{eq:small2}
\end{align}
Then for all $A$ with $\|A_e\|\le \delta$,
\begin{equation}\label{eq:Hess_lb_final}
\nabla^2 H_{j+1}(A)\ \succeq\ \varrho\,I,
\qquad
\varrho\ :=\ \alpha_{\min}\beta\cdot \frac{m(r_\ast;G)\,\lambda_\ast}{8}.
\end{equation}
In particular $\varrho>0$ is uniform in scale $j$ provided $\beta\ge \beta_{\min}>0$ and $\delta$ is fixed.
\end{theorem}

\begin{proof}[Structured proof with explicit constants]
Differentiate \eqref{eq:H_def}. By Lemma~\ref{lem:haar_convex}, the Haar Jacobian term is convex and can be discarded in a lower bound.
For each plaquette term $V_{\beta,\alpha_p}\circ F_p$, apply the chain rule:
\[
\nabla^2(V\circ F)=J^\ast (\nabla^2 V)J + \sum_{a} (\partial_a V)\,\nabla^2 F_a,
\]
where $J=\nabla F$ and $F_a$ are coordinates of $F$.
Using Lemma~\ref{lem:radial_hess} with $r_\ast$ and $\alpha_p\ge \alpha_{\min}$ gives
\[
J^\ast(\nabla^2 V)J\ \succeq\ \alpha_{\min}\beta\,m(r_\ast;G)\,J^\ast J.
\]
Summing over plaquettes yields the same bound with the block Jacobian $J(A)$ mapping $A$ to the vector $(F_p(A))_{p\in P(\mathcal Q)}$.

By Lemma~\ref{lem:BCH_bounds}, $J(A)=D_{\mathrm{red}}+\nabla N(A)$.
Hence, for any $v$,
\[
\|J(A)v\|\ \ge\ \|D_{\mathrm{red}}v\|-\|\nabla N(A)\|\|v\|
\ \ge\ (\sqrt{\lambda_\ast}-\varepsilon_1(\delta))\,\|v\|
\ \ge\ \tfrac12\sqrt{\lambda_\ast}\,\|v\|,
\]
using \eqref{eq:small1}. Therefore $J^\ast J\succeq (\lambda_\ast/4)I$.

For the second (nonlinear) term, use the gradient bound \eqref{eq:V_grad_bd}:
on $\|X\|\le r_\ast$,
\[
\big\|\sum_a (\partial_a V)\,\nabla^2 F_a\big\|_{\mathrm{op}}
\ \le\ \alpha_{\max}\beta\,M(r_\ast;G)\,r_\ast\,\varepsilon_2(\delta)
\ \le\ \beta\,M(r_\ast;G)\,r_\ast\,\varepsilon_2(\delta).
\]
With \eqref{eq:small2} this error is bounded by $(\alpha_{\min}\beta\,m(r_\ast;G)\lambda_\ast)/8$.
Combining this with the good term's lower bound $(\alpha_{\min}\beta\,m(r_\ast;G)\lambda_\ast)/4$ yields \eqref{eq:Hess_lb_final} with the stated factor $1/8$.
\end{proof}

\paragraph{Corollary ( convexity hypothesis).}
Theorem~\ref{thm:uniform_convexity} verifies Proposition~\ref{ass:uniform_convexity} with $\varrho$ given by \eqref{eq:Hess_lb_final}.

\subsection{Extracted-irrelevant scaling gap}

We now turn to the remaining finite-dimensional input needed for the irrelevant-sector contraction mechanism:
an explicit extracted-irrelevant scaling gap for the linearized RG map.

\begin{definition}[Engineering dimension and extraction level]\label{def:dim_extract}
Assign engineering dimensions in $d=4$:
\[
[A]=1,\qquad [F]=2,\qquad [D]=1.
\]
The relevant/marginal gauge-invariant local monomials have engineering dimension $\le 4$ and are spanned (in the gauge-invariant sector) by:
\[
\mathbf 1,\qquad \mathrm{tr}(F_{\mu\nu}F_{\mu\nu}),\qquad \mathrm{tr}(F_{\mu\nu}\widetilde F_{\mu\nu}).
\]
(Any scheme-specific additional marginal coordinates can be included here if needed.)
Define the extraction operator $\mathsf{Ext}$ to remove all local components of dimension $\le 4$ from polymer activities.
\end{definition}

\begin{lemma}[Scaling gap after extraction]\label{lem:Delta_min}
Assume extraction removes all components of engineering dimension $\le 4$ as in Definition~\ref{def:dim_extract}.
Then every remaining local monomial has dimension $\ge 6$.
Consequently the minimal dyadic scaling gap is
\[
\Delta_{\min}\ =\ 6-4\ =\ 2,
\]
and under deterministic dyadic rescaling (block size $L\mapsto 2L$),
\[
\|\mathrm{Rescale}(K)\|_{j+1}\ \le\ 2^{-\Delta_{\min}}\,\|K\|_j\ =\ \frac14\,\|K\|_j
\]
for extracted irrelevants, provided the seminorm $\|\cdot\|_j$ is defined with canonical scaling weights.
\end{lemma}

\begin{proof}
In $d=4$ pure gauge theory, there are no gauge-invariant monomials of odd dimension.
After removing dimension $\le 4$ components, the next possible gauge-invariant local operators are dimension $6$ (e.g.\ $\mathrm{tr}(F^3)$ and $\mathrm{tr}((D_\mu F_{\mu\nu})^2)$).
Dyadic rescaling reduces a dimension-$4+\Delta$ term by a factor $2^{-\Delta}$.
\end{proof}

\paragraph{Corollary ( scaling-gap hypothesis).}
Lemma~\ref{lem:Delta_min} verifies Proposition~\ref{ass:scaling_gap} with $\Delta_{\min}=2$.
Combined with the integration boundedness constant $C_E$ from the one-step integration estimate (Lemma~\ref{lem:E_bounded}),
the linear contraction constant is
\[
\kappa \ =\ C_E\,2^{-\Delta_{\min}}\ =\ \frac{C_E}{4},
\]
and the Lane~C stability mechanism requires $\kappa<1$.
